\documentclass[letterpaper,11pt]{article}
\usepackage[top=0.75in,bottom=0.75in,left=0.75in,right=0.75in]{geometry}
\usepackage{array}
\usepackage{longtable}
\usepackage{enumitem}
\usepackage{fancyhdr}
\usepackage{xcolor}
\usepackage[colorlinks=true,urlcolor=blue,citecolor=blue,linkcolor=blue]{hyperref}
\usepackage{xurl}

\urlstyle{same}
\setlength{\parindent}{0pt}
\setlength{\parskip}{4pt}
\pagestyle{fancy}
\fancyhf{}
\fancyfoot[L]{\footnotesize SCV Biographical Sketch v.2026-1}
\fancyfoot[R]{\footnotesize Page \thepage\ of \begin{NoHyper}\pageref{LastPage}\end{NoHyper}}
\renewcommand{\headrulewidth}{0pt}
\usepackage{lastpage}

\newcommand{\sechead}[1]{\vspace{6pt}\noindent\underline{\textbf{#1}}\vspace{2pt}\par}

\newcommand{\pubitemfont}{\normalsize}

\begin{document}

\noindent{\footnotesize Effective 06/16/2026}\hfill{\textbf{NASA BIOGRAPHICAL SKETCH}}\hfill{\footnotesize OMB-3145-0279}\par

\sechead{Identifying Information:}

\begin{tabular}{@{}p{6.5in}@{}}
\hline
\textbf{NAME:} Beck, Marcus W. \\
\hline
\textbf{PERSISTENT IDENTIFIER (PID):} \url{https://orcid.org/0000-0002-4996-0059} \\
\hline
\textbf{POSITION TITLE:} Senior Scientist \\
\hline
\textbf{PRIMARY ORGANIZATION AND LOCATION:} Tampa Bay Estuary Program, St. Petersburg, Florida, United States \\
\hline
\end{tabular}

\sechead{Professional Preparation:}

\renewcommand{\arraystretch}{1.5}
\begin{center}
\begin{tabular}{>{\centering\arraybackslash}p{2.3in}|>{\centering\arraybackslash}p{0.95in}|>{\centering\arraybackslash}p{1.1in}|>{\centering\arraybackslash}p{1.65in}}
\hline
\footnotesize ORGANIZATION AND LOCATION & \footnotesize DEGREE \newline\footnotesize(if applicable) & \footnotesize RECEIPT DATE & \footnotesize FIELD OF STUDY \tabularnewline
\hline
University of Minnesota, Twin Cities, Minnesota, United States & Ph.D. & 08/2009 - 06/2013 & Conservation Biology (Fisheries and Aquatic Biology); Statistics minor \tabularnewline
University of Minnesota, Twin Cities, Minnesota, United States & M.Sc. & 08/2007 - 06/2009 & Conservation Biology (Fisheries and Aquatic Biology) \tabularnewline
University of Florida, Gainesville, Florida, United States & B.Sc. & 08/2004 - 12/2006 & Zoology \tabularnewline
Santa Fe Community College, Gainesville, Florida, United States & A.A. & 08/2002 - 04/2004 & Zoology \tabularnewline
\end{tabular}
\end{center}
\renewcommand{\arraystretch}{1}

\sechead{Appointments and Positions}

\begin{longtable}{@{}p{0.9in}p{5.6in}@{}}
2019 - present & Senior Scientist, Tampa Bay Estuary Program, St. Petersburg, Florida, United States \\[4pt]
2021 - present & Private Contractor, EnviroDev, LLC, United States \\[4pt]
2017 - 2019 & Scientist, Southern California Coastal Water Research Project, Costa Mesa, California, United States \\[4pt]
2015 - 2017 & Post-Doctorate Researcher, Ecosystem Dynamics and Effects Branch, Gulf Ecology Division, United States Environmental Protection Agency, Gulf Breeze, Florida, United States \\[4pt]
2013 - 2015 & Post-Doctorate Research Fellow, Oak Ridge Institute for Science and Education, Ecosystem Dynamics and Effects Branch, Gulf Ecology Division, United States Environmental Protection Agency, Gulf Breeze, Florida, United States \\
\end{longtable}

\sechead{Products}

\begin{enumerate}[leftmargin=1.8em,itemsep=6pt,topsep=2pt]

\item Yang B, Liu Y, Law JA, Weisberg RH, D'Angelo S, Beck MW, Sherwood ET, Frazer TK. Tampa Bay Observing Network (TBON): A continuous data acquisition system in support of science and environmental management. Environmental Monitoring and Assessment. 2026; 198:1149. Available from: \url{https://doi.org/10.1007/s10661-026-16007-4}

\item Yang B, Liu Y, Tang W, Law J, Weisberg R, Beck MW, Sherwood ET, Frazer T. A machine learning model for real-time total nitrogen monitoring in Tampa Bay. Marine Pollution Bulletin. 2026; 233:120184. Available from: \url{https://doi.org/10.1016/j.marpolbul.2026.120184}

\item Beck MW, Sherwood ET, Murphy RR, Flaherty-Walia KE, Zhang Q, Simmons BA, Burke MC. Managing the unmanageable: Understanding the competing effects of salinity and nutrient loading on estuarine water quality. Estuaries and Coasts. 2026; 49(178):1-21. Available from: \url{https://doi.org/10.1007/s12237-026-01792-5}

\item Lizcano-Sandoval L, Beck MW, Scolaro S, Sherwood ET, Muller-Karger FE. Cloud-based satellite remote sensing for enhancing seagrass monitoring and ecosystem management. ISPRS Journal of Photogrammetry and Remote Sensing. 2025; 227:508-518. Available from: \url{https://doi.org/10.1016/j.isprsjprs.2025.06.034}

\item Arriola JM, Najjar RG, Brice\~{n}o H, Hu C, Herrmann M, Beck MW. Ecosystem metabolic rates estimated from diel oxygen measurements in two subtropical estuaries. Estuaries and Coasts. 2025; 48:155. Available from: \url{https://doi.org/10.1007/s12237-025-01597-y}

\item Medina M, Beck MW, Hecker J, Iadevaia N, Moody B, Anastasiou C, Tomasko D, Milbrandt EC, Kaplan D, Angelini C. Decoupled nutrient and chlorophyll trends in a large subtropical estuary: Water quality trends at the Greater Charlotte Harbor estuaries in southwest Florida. Estuaries and Coasts. 2025; 48:56. Available from: \url{https://doi.org/10.1007/s12237-025-01488-2}

\item Beck MW, Flaherty-Walia KE, Scolaro S, Burke MC, Furman BT, Karlen DJ, Pratt C, Anastasiou CJ, Sherwood ET. Hot and fresh: Evidence of climate-related suboptimal water conditions for seagrass in a large Gulf coast estuary. Estuaries and Coasts. 2024; 47:1475-1497. Available from: \url{https://doi.org/10.1007/s12237-024-01385-0}

\item Beck MW, Arriola JM, Herrmann M, Najjar RG. Fitting metabolic models to dissolved oxygen data: The Estuarine BAyesian Single-station Estimation method (EBASE). Limnology and Oceanography: Methods. 2024; 22(8):590-607. Available from: \url{https://doi.org/10.1002/lom3.10620}

\item Shi J, Hu C. South Florida estuaries are warming faster than global oceans. Environmental Research Letters. 2022 Dec 21; 18:014003. Available from: \url{https://doi.org/10.1088/1748-9326/aca8ba}

\item Beck MW, de Valpine P, Murphy R, Wren I, Chelsky A, Foley M, Senn DB. Multi-scale trend analysis of water quality using error propagation of generalized additive models. Science of the Total Environment. 2022; 802:149927. Available from: \url{https://doi.org/10.1016/j.scitotenv.2021.149927}

\item Beck MW, Altieri A, Angelini C, Burke MC, Chen J, Chin DW, Gardiner J, Hu C, Hubbard KA, Liu Y, Lopez C, Medina M, Morrison E, Phlips EJ, Raulerson GE, Scolaro S, Sherwood ET, Tomasko D, Weisberg R, Whalen J. Initial estuarine response to the nutrient-rich Piney Point release into Tampa Bay, Florida. Marine Pollution Bulletin. 2022; 178:113598. Available from: \url{https://doi.org/10.1016/j.marpolbul.2022.113598}

\item Beck MW, Burke MC, Raulerson GE, Scolaro S, Sherwood ET, Whalen J. Coordinated monitoring of the Piney Point wastewater discharge into Tampa Bay: Data synthesis and reporting. Florida Scientist. 2023; 86(2):288-300.

\item Beck MW, Wetherill B, Carr J. MassWateR: Improving quality control, analysis, and sharing of water quality data. PLOS ONE. 2023; 18(11):e0293737. Available from: \url{https://doi.org/10.1371/journal.pone.0293737}

\item Beck MW, Hagy JD, Le C. Quantifying seagrass light requirements using an algorithm to spatially resolve depth of colonization. Estuaries and Coasts. 2018; 41(2):592-610. Available from: \url{https://doi.org/10.1007/s12237-017-0287-1}

\item Le C, Lehrter JC, Hu C, MacIntyre H, Beck MW. Satellite observation of particulate organic carbon dynamics in two river-dominated estuaries. Journal of Geophysical Research: Oceans. 2017; 122(1):555-569. Available from: \url{https://doi.org/10.1002/2016JC012275}

\item Le C, Wu S, Hu C, Beck MW, Yang X. Phytoplankton decline in the eastern North Pacific transition zone associated with atmospheric blocking. Global Change Biology. 2019; 25(10):3485-3493. Available from: \url{https://doi.org/10.1111/gcb.14737}

\item Le C, Lehrter JC, Schaeffer BA, Hu C, Murrell MC, Hagy JD, Greene RM, Beck MW. Bio-optical water quality dynamics observed from MERIS in Pensacola Bay, Florida. Estuarine, Coastal and Shelf Science. 2016; 173:26-38. Available from: \url{https://doi.org/10.1016/j.ecss.2016.02.003}

\item Beck MW, Murphy R. Numerical and qualitative contrasts of two statistical models for water quality change in tidal waters. Journal of the American Water Resources Association. 2017; 53(1):197-219. Available from: \url{https://doi.org/10.1111/1752-1688.12489}

\item Beck MW, Schrandt MN, Wessel MR, Sherwood ET, Raulerson GE, Prasad AAB, Best BD. tbeptools: An R package for synthesizing estuarine data for environmental research. Journal of Open Source Software. 2021; 6(65):3485. Available from: \url{https://doi.org/10.21105/joss.03485}

\item Beck MW, de Valpine P, Murphy R, Wren I, Chelsky A, Foley M, Senn DB. wqtrends: Assess water quality trends with generalized additive models, v1.5.2. CRAN R package. 2026. Available from: \url{https://tbep-tech.github.io/wqtrends/}

\end{enumerate}

\sechead{Certification:}

I certify that the information provided is current, accurate, and complete. This includes but is not limited to information related to domestic and foreign appointments and positions.

I also certify that, at the time of submission, I am not a party to a malign foreign talent recruitment program.

In accordance with section 10634 of the CHIPS and Science Act of 2022 (42 U.S.C. \S\ 19234), I certify that I have completed the requisite research security training that meets the requirements specified in NASA's research security training requirements guidance within 12 months prior to proposal submission.

Misrepresentations and/or omissions may be subject to prosecution and liability pursuant to, but not limited to, 18 U.S.C. \S\S\ 287, 1001, 1031 and 31 U.S.C. \S\S\ 3729-3733 and 3802.

Certified by Beck, Marcus on 2026-10-08 18:02:31

\end{document}
